\subsubsection*{The finite condition after projecting to a character}
\label{subsec:ramified-character}

A finite class lifted through the crystalline lattices must still satisfy the required condition after projection to a character. The extension class below isolates the difference between an integral lift of a generator and a Galois-invariant one.

Let \(F_s=F_0(\mu_{\ell^s})\) be the full cyclotomic field of
Section~\ref{subsec:ramified-ht}, and let \(O\) be the coefficient DVR.
Choose a crystalline lattice extension and a lift \(v\) of its quotient generator:
\[
 0\longrightarrow T\longrightarrow\widetilde T\xrightarrow{p}O\longrightarrow0,\qquad p(v)=1.
\]
Put \(f(g)=gv-v\in T\).
Suppose that \(\psi:T\to E/O(-e)\) is an equivariant
map, where \(E\) is the fraction field of the
coefficient integers \(O\), and \(0\le e<t\le h\).
The kernel jumps lie in a common set
\(U\subseteq[-h,h]\) by strictness of the crystalline
filtration. We bound the loss to make
\(\psi f\) a finite class. Use \(M\) from
Section~\ref{subsec:ramified-ht} with the interval
\([-h,h]\).

\paragraph{The truncated integral periods.}
Write \(A_{\inf}\) for the integral period ring,
\[
 \xi=([\epsilon]-1)/\varphi^{-1}([\epsilon]-1)
\]
for the cyclotomic generator of the kernel of
Fontaine's map, and \(t_{\rm dR}\) for the
cyclotomic period.
Modulo \(\xi^{e+1}\), the cocycle
\(f\otimes1\) is a rational coboundary.
In fact strictness gives a lift of one in filtration
zero of the de Rham invariants of \(\widetilde T\).
Subtract it from \(v\) in the de Rham tensor.
The result lies in filtration zero of the tensor
with \(T[1/\ell]\), and has coboundary \(f\otimes1\).
The quotient \(A_{\inf}/\xi^{e+1}\) after inverting
\(\ell\) is the corresponding de Rham jet, so some
\(\ell\)-power gives an integral cobounding vector.

At level \(M\), the degree \(i\) piece of the jet
is the integral scalar ball in twist \(i\).
For Fontaine's map \(\vartheta\), the ratio of the two generators satisfies
\[
 \vartheta(t_{\rm dR}/\xi)=\zeta_\ell-1.
\]
The difference between the two
generators in degree \(i\) rescales the ball by an
element fixed by \(G_{F_M}\). Use the unprojected
frame \(D\) obtained at absolute depth \(s+n_0\) in
\eqref{eq:full-level-fontaine-map}. The map \(\vartheta\) in the jet
formula is the original depth-zero quotient. Indeed, writing
\(\mu=[\epsilon]-1\), one has
\[
  t_{\rm dR}/\xi
   =\bigl(\log(1+\mu)/\mu\bigr)\varphi^{-1}(\mu)
   \ \xrightarrow{\ \vartheta\ }\ 1\cdot(\zeta_\ell-1).
\]

Both the scalar in this formula and the scalar \(c\) in
\eqref{eq:full-level-fontaine-map} are fixed by \(G_{F_M}\).
The transformed lattice lies between the ball and
its \(\ell^{-1}\) multiple. Multiply a torsion cocycle
once by \(\ell\) to enter the ball. A rational
cobounding vector can also be brought into the ball
by a sufficiently large \(\ell\)-power.
The degree-one assertion of Lemma~\ref{lem:cyclotomic-invariant-loss}
applies even when the auxiliary extension has a zero jump. It bounds
the additional loss by
\[
   1+\max\left(1,\sum_{a\in U,\ a\ne i}(M+v_\ell(a-i))\right).
\]
This gives the summand of \(C(e+1)\) in degree \(i\)
in the original integral tensor lattice.

These losses add through the truncation.
For successive rational jets \(J'\twoheadrightarrow J\), the necessary lifting statement is
\[
 H^0(G_{F_M},J')\longrightarrow H^0(G_{F_M},J)\longrightarrow0 .
\]
An adapted filtration basis in the de Rham comparison reduces surjectivity to scalar filtered periods. In a scalar interval the
invariants come from its degree-zero piece, if
present, by Tate's vanishing in nonzero degrees.
This piece lifts by constants in the field \(F_M\).
After clearing denominators this is an integral
lift up to a power of \(\ell\).

For two successive quotients, kill the projected torsion class and subtract its integral coboundary. The remaining graded cocycle is torsion: some power of \(\ell\) gives a cobounding vector in the whole jet. Its
projection is invariant; lift that invariant
after clearing denominators and subtract the lift.
The remaining cobounding vector belongs to the
graded piece. If the projected class is killed by \(\ell^u\) and the recovered graded torsion by \(\ell^v\), the original class \(c\) satisfies
\[
 \ell^v(\ell^u c)=0.
\]
This is the additivity of the successive exponent bounds. Restriction to \(F_M\) and
corestriction back add
\[
 M-s=v_\ell([F_M:F_s])
\]
once. Thus \(\ell^{C(e+1)}(f\otimes1)\) is a
coboundary over \(A_{\inf}/\xi^{e+1}\).
\label{cor:cyclotomic-character-lifting-1}

\begin{lemma}[Integral dual exponentials of positive twists at two]
\label{lem:dyadic-integral-period}\label{hyp:positive-character-integral-period}
Let \(F_0/\mathbb Q_2\) be finite unramified, \(F_s=F_0(\mu_{2^s})\), \(s\ge2\), and \(e\ge1\). Identify the de Rham line for \(\mathbb Q_2(-e)\) with \(F_s\) in the Tate basis \(t_{\rm dR}^e\). The integral cohomology classes have dual exponentials satisfying
\[
  2^{1-e-v_2((e-1)!)}\mathcal O_{F_s}
           \subseteq\exp^*\!\bigl(H^1(F_s,\mathbb Z_2(-e))\bigr).
\]
\end{lemma}

\begin{proof}
\proofstep{Construct integral classes before taking their moments.}
Choose compatible primitive roots \(\zeta_r\) of order \(2^r\). On \(\mathcal O_{F_0}[[X]]\) write \(Y=1+X\) and let \(\varphi\) act by arithmetic Frobenius \(\sigma\) on coefficients and by \(Y\mapsto Y^2\). Its left inverse \(\psi\) extracts the even powers in the decomposition \(f=\varphi(f_0)+Y\varphi(f_1)\), so \(\psi(f)=f_0\).

At level \(r\ge2\), put \(u=1+2^r\). For a positive odd integer \(a\) and \(b\in\mathcal O_{F_0}\) set
\[
    g=b(Y^a-Y^{au}),\qquad y=\sum_{n\ge0}\varphi^n(g).
\]
Since \(Y^{2^n}-1\in(2,X)^{n+1}\), the series converges in \(\mathcal O_{F_0}[[X]]\). Both exponents in \(g\) are odd, hence \(\psi(g)=0\) and \(\psi(y)=y\). The integral Iwasawa comparison for the trivial lattice attaches to \(y\) a system of integral cohomology classes. Its moment of order \(-e\) at level \(r\) belongs to \(H^1(F_r,\mathbb Z_2(-e))\). At two the torsion-free cyclotomic tower used for this comparison starts at \(r=2\). \cite[Proposition~II.1.2 and Theorem~II.1.3]{cherbonnier-colmez-iwasawa}

\proofstep{Evaluate a finite-level moment.}
Let \(y^{\sigma^{-r}}\) denote the series obtained by acting on its coefficients. In the chosen Tate basis, the dual exponential of the moment is the coefficient of \(t^e\) in
\[
   2^{-r}y^{\sigma^{-r}}(\zeta_r\exp(t/2^r)-1),\qquad t=t_{\rm dR}.
\]
The dual exponential formula of Cherbonnier and Colmez gives this series at all sufficiently high levels. To descend one step from \(j\) to \(j-1\ge2\), trace the series. The two conjugates of \(\zeta_j\) have opposite signs. Their average on \(y\) retains its even part \(\varphi(\psi y)=\varphi(y)\). Hence
\[
 \operatorname{Tr}_{F_j/F_{j-1}}\!
     \left(2^{-j}y^{\sigma^{-j}}(\zeta_j\exp(t/2^j)-1)\right)
 =2^{-(j-1)}y^{\sigma^{-(j-1)}}(\zeta_{j-1}\exp(t/2^{j-1})-1).
\]
Corestriction of the moment commutes with the dual exponential and this unnormalized trace. The coefficient formula therefore holds for each \(r\ge2\). \cite[Theorem~IV.2.1 and its proof]{cherbonnier-colmez-iwasawa}

After division by the odd unit \(a^e\), the computed value is \(C_r S_{r,a}(b)\), where
\[
 C_r=\frac{1-u^e}{2^{r(e+1)}e!},\qquad
 S_{r,a}(b)=\sum_{n\ge0}2^{ne}\sigma^{n-r}(b)\zeta_r^{a2^n}.
\]
Indeed the powers \(a2^n\) and \(au2^n\) give the same root of unity at level \(r\), while the coefficient of \(t^e\) in their exponentials introduces their \(e\)-th powers. The series for this coefficient converges since \(e>0\). For \(r\ge2\), the lifting-the-exponent formula yields
\[
    v_2(C_r)=r+v_2(e)-r(e+1)-v_2(e!)
            =-re-v_2((e-1)!).
\]

\proofstep{Recover each integral direction in the cyclotomic field.}
Replace \(a\) by \(a+2^{r-1}\). In the difference of the two normalized moments all the terms with \(n\ge1\) cancel. Since \(b\) is arbitrary,
\[
 C_r\bigl(S_{r,a}(b)-S_{r,a+2^{r-1}}(b)\bigr)
                   =2 C_r\sigma^{-r}(b)\zeta_r^a .
\]
Thus the image at level \(r\) contains \(2^{1-re-v_2((e-1)!)}\mathcal O_{F_0}\zeta_r^a\) for each odd \(a\).

The ground-field direction is also needed. At \(r=2\), add the normalized moments for \(a=1,3\). Their leading terms cancel, the roots in the next terms equal \(-1\), and all later roots equal one. Consequently
\[
 C_2\bigl(S_{2,1}(b)+S_{2,3}(b)\bigr)
 =2^{e+1}C_2\left(-\sigma^{-1}(b)+
             \sum_{n\ge2}2^{(n-1)e}\sigma^{n-2}(b)\right).
\]
The endomorphism in parentheses is congruent to \(-\sigma^{-1}\) modulo \(2\), hence is bijective on the finite free \(\mathbb Z_2\)-module \(\mathcal O_{F_0}\). The image contains \(2^{1-e-v_2((e-1)!)}\mathcal O_{F_0}\).

Restriction from \(F_r\) to \(F_s\) for \(2\le r\le s\) preserves integrality and includes the dual exponential value into \(F_s\). Each power of \(\zeta_s\) is \(\pm1\) or an odd power of a primitive \(\zeta_r\) with \(2\le r\le s\). These powers span \(\mathcal O_{F_0}[\zeta_s]=\mathcal O_{F_s}\) over \(\mathcal O_{F_0}\). The exponent \(1-re-v_2((e-1)!)\) is at most \(1-e-v_2((e-1)!)\), proving the assertion.
\end{proof}

\begin{proposition}[Returning scalar cocycles to the finite condition]\label{prop:cyclotomic-scalar-finite-condition}
Take the cyclotomic module of Proposition~\ref{prop:ramified-integral-wach-comparison} for the ramified lattice. Let \(f(g)=gv-v\) represent the crystalline lattice extension over \(F_s\),
and let \(\psi:T\to(E/O)(-e)\) be equivariant, with \(0\le e<t\le h\).
With the integer loss \(C(k)\) of \eqref{eq:local-jump-loss}, put
\[
 A_e=(e+2)s+v_\ell(e)+\lfloor e/(\ell-1)\rfloor+1,
 \qquad B_e=C(e+1)+s+v_\ell(e+1)+A_e,
\]
where the \(v_\ell(e)\) term in \(A_e\) is omitted when \(e=0\). Then the pushed-out class satisfies
\[
 \ell^{B_e}[\psi\circ f]\in H^1_f(F_s,(E/O)(-e)).
\]
Here the finite subgroup is the image of the rational finite condition.
For a class defined over the original local field at level \(N\),
corestriction adds at most \(1+\lfloor\log_\ell N\rfloor\) to the exponent.
\end{proposition}

\begin{proof}
\proofstep{Use the Bockstein allowance to lift the divisible cocycle to rational coefficients.}
Lift \(\psi\) on an integral
\(\mathbb Z_\ell\)-basis of \(T\) to the rational
target. The lift need only be equivariant modulo the
target lattice. Expand in an integral
\(\mathbb Z_\ell\)-basis of this target lattice.
Each coordinate of \(\psi f\) modulo \(\mathbb Z_\ell\)
is now a cocycle in
\(\mathbb Q_\ell/\mathbb Z_\ell(-e)\).

By local Tate duality,
\(H^2(F_s,\mathbb Z_\ell(-e))\) is dual to the
invariants in \(\mathbb Q_\ell/\mathbb Z_\ell(e+1)\).
The cyclotomic generator has \(v_\ell(\chi(\gamma)^{e+1}-1)=s+v_\ell(e+1)\). Duality therefore gives
\[
 \ell^{s+v_\ell(e+1)}H^2(F_s,\mathbb Z_\ell(-e))=0.
\]
The integral Bockstein sequence
then gives a rational cocycle representing the scaled
divisible cocycle in each coordinate. Adjust by a
rational coboundary lifting the discrepancy between
the divisible cocycles. The representatives then
agree pointwise modulo
\(\mathbb Z_\ell\).
The equivariance errors of the lifted coefficient
map are integral, so applying it to the jet
coboundary preserves the congruence modulo integral
periods. After scaling further by \(\ell^{C(e+1)}\),
each rational cocycle \(z\) so obtained satisfies
\[
 z(g)\equiv \chi(g)^{-e}g b-b
       \pmod {A_{\inf}+t_{\rm dR}^{e+1}B_{\rm dR}^+}
 \qquad (b\in B_{\rm dR}^+).
\]
The congruence means that the error is in the
image of \(A_{\inf}\) in the jet.

Use the dual exponential in the Tate basis \(t_{\rm dR}^e\). Its value \(\alpha\in F_s\) is
specified by the remaining term
\(\log\chi(g)\alpha t_{\rm dR}^e\) after a rational
twisted coboundary in that jet. To see this, filter the rational jet
\((B_{\rm dR}^+/t_{\rm dR}^{e+1})(-e)\) by powers of \(t_{\rm dR}\).
Its degree-\(i\) quotient is \(\mathbb C_\ell(i-e)\) for \(0\le i\le e\).
For \(i<e\), Tate's theorem gives
\(H^1(F_s,\mathbb C_\ell(i-e))=0\), so the cocycle in that quotient is
a coboundary. In degree \(e\) the twist is trivial, and Tate's calculation gives
\[
 H^1(F_s,\mathbb C_\ell)=F_s\,[\log\chi].
\]

When the dual exponential is zero, the Bloch--Kato
comparison gives a cobounding vector over
\(B_{\mathrm{cris}}\otimes_{\mathbb Q_\ell}K_\ell(-e)\),
where \(K_\ell\) is the coefficient field. This identifies
the corrected class with the crystalline finite condition.
\cite{tate-pdivisible,bloch-kato}\label{cor:cyclotomic-character-lifting-2}

\proofstep{Estimate the dual exponential in the truncated period coordinates.}
The Bockstein argument leaves a rational lift. Control of its projection to the finite subgroup now comes from the integral periods entering the dual exponential.

The displayed integral congruence bounds \(\alpha\).
At degree \(i<e\) use the integral multiplier
\[
 \ell^{1+s+v_\ell(i-e)}
\]
on the modified cobounding vector, then subtract an integral lift of its leading coefficient.
The degree \(i\) scalar belongs to the ball
modulo a twisted invariant, after this multiplication.
This is the rank-one ball argument at depth \(s\)
in twist \(i-e\), using the valuation of
\(\chi(\gamma)^{i-e}-1\).
In the \(t_{\rm dR}^i\) coordinates the integral
leading coefficient is multiplied by
\((\zeta_\ell-1)^i\), from \(t_{\rm dR}/\xi\).
Thus it can be lifted by Fontaine's surjection
and the power \(\xi^i\).
After eliminating all lower degrees the logarithm coefficient is
\[
 \ell^{e(s+1)+v_\ell(e!)}(\zeta_\ell-1)^e\alpha.
\]
Its two factors record the scalar elimination costs and the generator conversion, respectively.

If \(\alpha_1\log\chi(g)+g b_1-b_1\) is integral
and \(\alpha_1\in F_s\), apply the integral
coboundary estimate on the cyclotomic kernel
after multiplying by \(\ell\).
The remaining cobounding vector is fixed by that
kernel. Take the integral normalized trace at the
generator. It kills the difference of this
vector and its translate and gives
\(\ell\alpha_1\log\chi(\gamma)\) integral.
Since \(v_\ell(\log\chi(\gamma))=s\), one obtains
\[
 v_\ell(\alpha)\ge
 -e(s+1)-v_\ell(e!)-1-s-\frac{e}{\ell-1}.
\]

At weight zero the relevant integral image can be computed by local class field theory. Write \(E_0=(\ell-1)\ell^{s-1}\) for the ramification index of \(F_s/F_0\). The different is generated by \(\Phi'_{\ell^s}(\zeta_{\ell^s})\); its absolute valuation, with \(v_\ell(\ell)=1\), is \(s-1/(\ell-1)\). For \(v_\ell(z)\ge1/E_0\) the logarithmic series gives
\[
 v_\ell(\log(1+z))\ge \min_{n\ge1}\{n/E_0-v_\ell(n)\}
                  =1+1/(\ell-1)-s .
\]
For fixed \(v_\ell(n)=j\) the minimum occurs at \(n=\ell^j\); comparing consecutive values gives the equality. For the inverse different \(\mathfrak D^{-1}=\mathfrak D_{F_s/\mathbb Q_\ell}^{-1}\) this gives
\[
          \ell^{-1}\mathcal O_{F_s}\log(1+\mathfrak m_{F_s})
                              \subseteq\mathfrak D^{-1}.
\]

For \(a\in\ell^{-1}\mathcal O_{F_s}\) define on principal units
\[
       h_a(u)=\operatorname{Tr}_{F_s/\mathbb Q_\ell}(a\log u)\in\mathbb Z_\ell .
\]
Extend it by zero on prime-to-\(\ell\) roots of unity and a uniformizer. Local reciprocity gives an integral Galois character. The duality of the Bloch--Kato exponential with Kummer theory identifies its dual exponential with \(a\), up to the common reciprocity sign: pairing on principal units on both sides is the displayed trace. This proves the required inclusion for \(e=0\). \cite[Section~3]{bloch-kato-finite}

For odd \(p=\ell\) and \(e>0\), put \(m=e+1\) and \(L=F_s\). For \(m\ge2\) the rational Bloch--Kato exponential is an isomorphism. In its standard Tate basis put
\[
 \mathcal H_m=\exp_{L,m}^{-1}\!\left(\operatorname{im}\big[H^1(L,\mathbb Z_p(m))
                           \longrightarrow H^1(L,\mathbb Q_p(m))\big]\right)\subset L .
\] Let \(\sigma\) be Frobenius on \(F_0\), and let \(f\in\mathcal O_{F_0}[[X]]\) range over the power series satisfying \(\sum_{\zeta^p=1}f(\zeta(1+X)-1)=0\). Denote their module by \(R\).

The integral Coleman sequence and its explicit evaluation at level \(s\) give
\[
          p^{s+v_p(e)}\mathcal H_m\ \subseteq\ e!\,\Xi_{m,s}(R).
\]
Indeed the quotient of the full inverse-limit classes by the classes from principal units, after passing to the coinvariants for \(F_s\), is \(H^0(F_s,\mathbb Q_p/\mathbb Z_p(e))\). A generator with cyclotomic value \(1+p^s\) acts on that quotient through the relation
\[
       ((1+p^s)^e-1)x=0,\qquad v_p((1+p^s)^e-1)=s+v_p(e).
\] For the principal-unit classes the evaluation diagram sends each integral class to \((m-1)!\,\Xi_{m,s}(f)\) in exponential coordinates. Here \(L=F_s\) is already the entire cyclotomic level, so the trace from that level to \(L\) in the diagram is the identity. \cite[Sections~2.3.1--2.3.2, (4)--(5)]{benois-nguyen-tamagawa}

The explicit evaluation gives more than its determinant. In the formula for \(\Xi_{m,s}\), division by \(p^{m-s}\) makes the terms of index \(0\le k<s-1\) divisible by \(p^m\): their coefficients are \(p^{m(s-1-k)}\). The term for \(k=s-1\) and the constant correction give
\[
 p^{s-m}\Xi_{m,s}(f)\equiv
 f^{\sigma^{-1}}(\zeta_p-1)-(1-p^m\sigma^{-1})^{-1}f^{\sigma^{-1}}(0)
                      \pmod{p^m\mathcal O_L}.
\]
The inverse in the second term is the convergent geometric series and is the identity modulo \(p^m\). Evaluation of an integral series at \(\zeta_p-1\) agrees with its constant coefficient modulo \(\zeta_p-1\). The calculation also covers \(s=1\), when the earlier sum is empty. \cite[Section~2.2.1, definition of \(\Xi_{m,n}\)]{benois-nguyen-tamagawa}

In particular every vector in \(\mathcal H_m\) has valuation at least
\[
 A=v_p(e!)+m-2s-v_p(e)+1/(p-1).
\]
Under local Tate duality the image of \(H^1(L,\mathbb Z_p(1-m))\) in the dual de Rham line is the trace dual of \(\mathcal H_m\). Indeed the integral cup product on the torsion-free quotients is perfect, and the Bloch--Kato exponential and dual exponential are adjoint for this pairing. In the dual Tate bases their de Rham pairing is \((a,b)\mapsto\operatorname{Tr}_{L/\mathbb Q_p}(ab)\). \cite[Section~3]{bloch-kato-finite}

Put \(\beta=s-1-e-v_p((e-1)!)\). The valuation threshold for the inverse different is \(-s+1/(p-1)\); with \(m=e+1\) one has
\[
        \beta+A=-s+1/(p-1).
\]
Thus \(p^\beta\mathcal O_L\mathcal H_m\subseteq\mathfrak D^{-1}\). Every vector of \(p^\beta\mathcal O_L\) has integral trace pairing with \(\mathcal H_m\), proving the required inclusion at every odd prime.

For \(\ell=2\) and \(e>0\), the exponent in Lemma~\ref{lem:dyadic-integral-period} is independent of the cyclotomic level. Since \(s\ge2\), the required ball satisfies
\[
  2^{s-1-e-v_2((e-1)!)}\mathcal O_{F_s}
     \subseteq 2^{1-e-v_2((e-1)!)}\mathcal O_{F_s}
     \subseteq\exp^* H^1(F_s,\mathbb Z_2(-e)).
\]
\label{cor:cyclotomic-character-lifting-3}

Set
\[
 A_e=(e+2)s+v_\ell(e)+\lfloor e/(\ell-1)\rfloor+1,
\]
omitting the \(v_\ell(e)\) term at \(e=0\).
The bound for \(\alpha\) gives
\(\ell^{A_e}\alpha\in\ell^\beta\mathcal O_{F_s}\).
For positive \(e\) use
\(v_\ell(e!)=v_\ell(e)+v_\ell((e-1)!)\) and
\(e/(\ell-1)<\lfloor e/(\ell-1)\rfloor+1\).
For \(e=0\) it follows directly.
Subtract an integral cocycle with this exponential
from \(\ell^{A_e}z\). The difference has zero
exponential and represents the same divisible
class. Reassemble the scalar coordinates in the
coefficient lattice. This puts the character image
in the divisible finite condition after loss
\(C(e+1)+s+v_\ell(e+1)+A_e\).
Corestriction adds at most the original local degree
loss, giving the character allowance used above.
\end{proof}

